% GENERATED FILE. Edit canonical lessons, then run npm run generate:books.
% canonical-source-hash: fnv1a64:b7a2c8d75951282a
% canonical-lessons: FR-C352-empecher, FR-C352-entraver, FR-C352-favoriser, FR-C352-atout, FR-C352-frein

\chapter{Freins et atouts}
\label{ch:fr-freins-et-atouts}
\hlchaptermodality{352}

\begin{chapteropening}
\textbf{By the end of this chapter:} \emph{I can say what stands in the way of a plan, what prevents or hinders it, and what favours it or works as an asset.}

Complete the last lesson of chapter 352: I can say what stands in the way of a plan, what prevents or hinders it, and what favours it or works as an asset.
\end{chapteropening}

\section[empêcher]{empêcher — to prevent, to stop (from)}

\label{lesson:FR-C352-empecher}



\begin{quote}
Before the new one: say the French for to justify, then the French for the justification.
\end{quote}

\subsection*{You'll want to know: empêcher}
\textbf{empêcher} — \textquotedblleft{}to prevent, to stop (from)\textquotedblright{}.

\begin{quote}
Le bruit m'empêche de dormir.
\end{quote}

\begin{quote}
\emph{The noise stops me from sleeping.}
\end{quote}

\emph{Empêcher quelqu'un de faire quelque chose}. \emph{Il n'empêche que} — all the same.

\begin{cousinweb}
\emph{empêcher} is inherited through Old French \emph{empeechier} from Late Latin \emph{impedicāre} \textquotedblleft{}to fetter, to hinder\textquotedblright{} (EtymDB), built on \emph{pedica} \textquotedblleft{}a fetter for the foot\textquotedblright{} (Lewis \& Short).
\end{cousinweb}

\begin{grammarlens}[title={What you've built}]
One more word for telling, explaining and describing at B1.
\end{grammarlens}

\subsection*{Your turn}
\begin{itemize}
\raggedright
  \item \emph{Say it:} \emph{empêcher}
  \item \emph{Say it:} \emph{empêcher}, once more
  \item \emph{Say it:} a sentence of your own with \emph{empêcher}
  \item \emph{Recall:} say the French for to justify, then the French for the justification, then say \emph{empêcher} again
\end{itemize}

\begin{tcolorbox}[breakable,colback=teal!4,colframe=teal!35!black,title={Before you move on}]
What does empêcher mean? (\textquotedblleft{}To prevent, to stop (from)\textquotedblright{}.) Say it once more.
\end{tcolorbox}
\section[entraver]{entraver — to hinder, to hamper}

\label{lesson:FR-C352-entraver}



\begin{quote}
Before the new one: say the French for the justification, then the French for to prevent.
\end{quote}

\subsection*{You'll want to know: entraver}
\textbf{entraver} — \textquotedblleft{}to hinder, to hamper\textquotedblright{}.

\begin{quote}
La neige entrave le travail des pompiers.
\end{quote}

\begin{quote}
\emph{The snow is hampering the firefighters.}
\end{quote}

Formal. \emph{Une entrave} was a chain on an animal's legs.

\begin{cousinweb}
\emph{entraver} is usually explained as built from \emph{en-} and Old French \emph{tref} \textquotedblleft{}beam\textquotedblright{}, from Latin \emph{trabs} \textquotedblleft{}a beam\textquotedblright{} (Lewis \& Short): a piece of wood tied to the legs. The details are not settled.
\end{cousinweb}

\begin{grammarlens}[title={What you've built}]
One more word for telling, explaining and describing at B1.
\end{grammarlens}

\subsection*{Your turn}
\begin{itemize}
\raggedright
  \item \emph{Say it:} \emph{entraver}
  \item \emph{Say it:} \emph{entraver}, once more
  \item \emph{Say it:} a sentence of your own with \emph{entraver}
  \item \emph{Recall:} say the French for the justification, then the French for to prevent, then say \emph{entraver} again
\end{itemize}

\begin{tcolorbox}[breakable,colback=teal!4,colframe=teal!35!black,title={Before you move on}]
What does entraver mean? (\textquotedblleft{}To hinder, to hamper\textquotedblright{}.) Say it once more.
\end{tcolorbox}
\section[favoriser]{favoriser — to favour, to encourage}

\label{lesson:FR-C352-favoriser}



\begin{quote}
Before the new one: say the French for to prevent, then the French for to hinder.
\end{quote}

\subsection*{You'll want to know: favoriser}
\textbf{favoriser} — \textquotedblleft{}to favour, to encourage\textquotedblright{}.

\begin{quote}
Le sport favorise un bon sommeil.
\end{quote}

\begin{quote}
\emph{Exercise encourages good sleep.}
\end{quote}

To help something happen; also to treat someone better than others.

\begin{cousinweb}
\emph{favoriser} is built on \emph{faveur} \textquotedblleft{}favour\textquotedblright{}, from Latin \emph{favor} \textquotedblleft{}goodwill, favour\textquotedblright{} (Lewis \& Short), from \emph{favēre} \textquotedblleft{}to be well disposed\textquotedblright{}.
\end{cousinweb}

\begin{grammarlens}[title={What you've built}]
One more word for telling, explaining and describing at B1.
\end{grammarlens}

\subsection*{Your turn}
\begin{itemize}
\raggedright
  \item \emph{Say it:} \emph{favoriser}
  \item \emph{Say it:} \emph{favoriser}, once more
  \item \emph{Say it:} a sentence of your own with \emph{favoriser}
  \item \emph{Recall:} say the French for to prevent, then the French for to hinder, then say \emph{favoriser} again
\end{itemize}

\begin{tcolorbox}[breakable,colback=teal!4,colframe=teal!35!black,title={Before you move on}]
What does favoriser mean? (\textquotedblleft{}To favour, to encourage\textquotedblright{}.) Say it once more.
\end{tcolorbox}
\section[l'atout]{l'atout — the asset, the strong point; the trump card}

\label{lesson:FR-C352-atout}



\begin{quote}
Before the new one: say the French for to hinder, then the French for to favour.
\end{quote}

\subsection*{You'll want to know: l'atout}
\textbf{l'atout} — \textquotedblleft{}the asset, the strong point; the trump card\textquotedblright{}.

\begin{quote}
Parler trois langues est un vrai atout.
\end{quote}

\begin{quote}
\emph{Speaking three languages is a real asset.}
\end{quote}

In card games, \emph{l'atout} is the trump suit. Masculine.

\begin{cousinweb}
\emph{atout} is built from \emph{à} and \emph{tout} \textquotedblleft{}all\textquotedblright{}: a card that beats all the others. \emph{tout} is inherited from Latin \emph{tōtus} \textquotedblleft{}whole\textquotedblright{} (Lewis \& Short).
\end{cousinweb}

\begin{grammarlens}[title={What you've built}]
One more word for telling, explaining and describing at B1.
\end{grammarlens}

\subsection*{Your turn}
\begin{itemize}
\raggedright
  \item \emph{Say it:} \emph{l'atout}
  \item \emph{Say it:} \emph{l'atout}, once more
  \item \emph{Say it:} a sentence of your own with \emph{l'atout}
  \item \emph{Recall:} say the French for to hinder, then the French for to favour, then say \emph{l'atout} again
\end{itemize}

\begin{tcolorbox}[breakable,colback=teal!4,colframe=teal!35!black,title={Before you move on}]
What does l'atout mean? (\textquotedblleft{}The asset, the strong point; the trump card\textquotedblright{}.) Say it once more.
\end{tcolorbox}
\section[le frein]{le frein — the brake; the obstacle}

\label{lesson:FR-C352-frein}



\begin{quote}
Before the new one: say the French for to favour, then the French for the asset.
\end{quote}

\subsection*{You'll want to know: le frein}
\textbf{le frein} — \textquotedblleft{}the brake; the obstacle\textquotedblright{}.

\begin{quote}
Le prix est un frein pour beaucoup de familles.
\end{quote}

\begin{quote}
\emph{The price is an obstacle for many families.}
\end{quote}

On a car or a bike, \emph{les freins} are the brakes; the verb is \emph{freiner}.

\begin{cousinweb}
\emph{frein} is inherited from Latin \emph{frēnum} \textquotedblleft{}a bridle, a bit\textquotedblright{} (EtymDB; Lewis \& Short): what holds a horse back.
\end{cousinweb}

\begin{grammarlens}[title={What you've built}]
One more word for telling, explaining and describing at B1.
\end{grammarlens}

\subsection*{Your turn}
\begin{itemize}
\raggedright
  \item \emph{Say it:} \emph{le frein}
  \item \emph{Say it:} \emph{le frein}, once more
  \item \emph{Say it:} a sentence of your own with \emph{le frein}
  \item \emph{Recall:} say the French for to favour, then the French for the asset, then say \emph{le frein} again
\end{itemize}

\begin{tcolorbox}[breakable,colback=teal!4,colframe=teal!35!black,title={Before you move on}]
What does le frein mean? (\textquotedblleft{}The brake; the obstacle\textquotedblright{}.) Say it once more.
\end{tcolorbox}
